\section*{Ensayo por José Manuel Núñez Ruíz}
\addcontentsline{toc}{subsubsection}{Ensayo por José Manuel Núñez Ruíz.}

¿Alguna vez te has preguntado qué hace a Google o Meta empresas multimillonarias si sus servicios son ``gratuitos''? o ¿Por qué tenemos que aceptar cookies en todas partes? Estas son preguntas que desembocan en un análisis que tal vez no se hace la mayoría de la gente al aceptar cookies o firmar los términos y condiciones de la mayoría de aplicaciones que usa. Los datos se han vuelto el nuevo activo que las empresas codician; estos les permiten ganar terreno contra otras empresas competidoras. La industria se ha convertido en quién te vende mejor ciertos productos, y para ello necesitan recolectar la mayor cantidad de información para procesarla y convertirla en conocimiento. Estos datos pueden ser recolectados de todas partes: las redes sociales (Instagram, Facebook, X, etc.), de tu misma casa con dispositivos electrónicos conectados a internet (Alexa, electrodomésticos inteligentes, etc.) y, por supuesto, una de las más recientes, el uso masivo de Inteligencia Artificial.

A primera vista esto no parece tan malo, ya que la recompensa es obtener acceso a información de forma masiva. Actualmente, aplicaciones de geolocalización como Google Maps, los navegadores y diversas aplicaciones más nos permiten tener una vida más ``cómoda''; sin embargo, debemos concientizarnos sobre el uso de nuestros datos de forma no transparente. En lo personal, creo que un cambio radical en el sistema es algo improbable y no eficiente, pues aún hay sectores de la población muy grandes cuyas prioridades están muy lejos de esta cuestión; sin embargo, concientizarnos y tener otras opciones sobre el uso de nuestros datos y la transparencia de ello son los primeros pasos que se deben implementar en el sistema.

En el artículo (\cite{becerril2018value}) se habla primordialmente de que solemos quitarle valor a nuestros datos. Como se menciona en un estudio, el usuario promedio calificaba con el valor de una hamburguesa a sus datos personales; aunque esto puede sonar exagerado, la autora quiere transmitirnos que no valoramos lo suficiente nuestros datos. En efecto, empresas gigantes como Google o Meta obtienen gran parte de sus ganancias (más del 90\%) gracias a nosotros y a los datos que obtienen. ¿Cómo lo hacen? Venden la información recopilada a anunciantes para atraer más nuestra atención, la utilizan para predicciones del mercado o dejan que nosotros mismos mejoremos sus sistemas a través de opiniones. Esto último me hace recordar que cuando estaba en la preparatoria, mis amigos estaban extrañados de que cuando consultaban algo recientemente, digamos veían ``Stranger Things'', en su feed de redes sociales o aun en recomendaciones de Google se multiplicaba la publicidad, aun cuando la serie la habían visto en Netflix, que parecería no tener conexión con las redes sociales; esto, sin embargo, muestra cómo no estamos conscientes de la recopilación masiva de nuestros datos. Adicionalmente, se menciona la emersión del Internet de las Cosas (IoT) para recopilar aún más datos ahora en tu casa o en distintas partes de las que no somos conscientes. Este Big Data es además poseído por empresas privadas, las cuales no piensan poner en riesgo la competencia y no lo hacen transparente.

En el desarrollo del artículo, la autora, aunque menciona que no hay una forma exacta de calcular el valor de nuestros datos, cita distintos artículos donde se muestran ejemplos de cuánto podrían generar los usuarios con sus datos personales o respondiendo encuestas. En mi opinión, esta parte no es demasiado relevante (buscar un número para darte una idea), pues el valor de tus datos, y más de tus datos sensibles, dependen del contexto: ¿Hay un mal uso de ellos? ¿Los otorgas al registrarte en distintas páginas de internet? Y además, qué tan interesado estás en cuidarlos o si estás dispuesto a otorgarlos por algún beneficio. En mi práctica profesional, al diseñar una Base de Datos o ser un ABD, tengo la responsabilidad de no solo diseñar un sistema eficiente, sino de seguir normativas éticas y respetar la privacidad de los usuarios implementando estrictas medidas de seguridad.

Considero que una posición extremista en esto es inútil y no realista; pensar que de la noche a la mañana seremos los únicos dueños de todos nuestros datos y que podremos recibir lo que merecemos de estos es un camino que no va a pasar. A pesar de esto, apoyo los argumentos del artículo sobre la concientización de esto, apoyados con números sobre la desigualdad de poder, la gran cantidad de Big Data que producimos y la poca rentabilidad económica que recibimos. Podemos notar que nada de esto sería posible sin la existencia de las Bases de Datos. Aun ahora, con los conjuntos masivos de datos que existen (Big Data), las bases de datos nos permiten relacionar los datos para que sean útiles; no les serviría a las empresas tener cantidades exorbitantes de datos si no tienen forma de organizarlos y extraer información para posteriormente asimilar conocimiento de estos.

En conclusión, la vulnerabilidad de nuestros datos es, en efecto, algo de lo que debemos estar conscientes, más aún sabiendo que la mayor parte de la población no tiene idea de lo que acepta en los términos y condiciones o qué son las cookies. Sin embargo, la realidad de las redes sociales y la comodidad que nos ofrecen también es innegable, por lo que debemos encontrar un balance y luchar por la transparencia de nuestros datos sin caer en el extremismo de no aceptar ningunos términos y condiciones que se nos propongan. 

En mi caso, la autora logró hacerme más consciente de la importancia de valorar mis datos personales y ahora me detendré a leer más atentamente las cookies o los términos y condiciones; también me parece muy interesante la idea del Almacenamiento de Datos Personales (PDS), que otorga alternativas viables a la problemática planteada mediante tecnologías centradas en el usuario como dueños de su propia información.

A futuro quedan diversas incógnitas, pero la más relevante para mí sería: ¿cómo se controlará el manejo de datos con las IAs? Ya que la inteligencia artificial consume innumerables cantidades de información que además puede no tener integridad, así que regular su alimentación parece algo que debe plantearse próximamente, ya que los modelos de aprendizaje profundo suelen entrenarse mediante minería de datos masiva de fuentes públicas sin un consentimiento explícito ni mecanismos de control, lo que pone en riesgo directamente los principios de privacidad y gobernanza de datos que se intentan promover actualmente.
